\section{一 正方形内的旋转}
\begin{problem}{题 1 旋转拼出线段和}
正方形 \(ABCD\) 中，\(E\)、\(F\) 分别在边 \(BC\)、\(CD\) 的内部，\(\angle EAF=45^\circ \)。证明 \(EF=BE+DF\)。再思考：若只知道 \(EF=BE+DF\)，能否推出 \(\angle EAF=45^\circ \)？
\diagram{\figOne}
\hint{延长 \(CD\) 至 \(G\)，使 \(DG=BE\)，再比较 \(AE\) 与 \(AG\)。}
\end{problem}
\begin{problem}{题 2 四个条件相互推导}
正方形 \(ABCD\) 中，\(M\)、\(N\) 分别在 \(AD\)、\(CD\) 的内部。证明下列四个条件彼此等价：\((1) \angle MBN=45^\circ \)；\((2) MN=AM+CN\)；\((3) \angle AMB=\angle BMN\)；\((4) \angle BNM=\angle BNC\)。
\diagram{\figTwo}
\hint{先由\((1)\)证明\((2)\)。再从 \(B\) 向 \(MN\) 作垂线，观察等角与等距的联系。}
\end{problem}
\clearpage
\begin{problem}{题 3 从等角恢复半角}
正方形 \(ABCD\) 中，\(M\in AD\)、\(N\in CD\)，\(\angle BMN=\angle MBC\)，\(AB=9\)，\(DM=6\)，求 \(CN\)。
\diagram{\figThree}
\hint{由 \(AD\parallel BC\)，把已知等角转写为题 \(2\) 中的一个条件。}
\end{problem}
\begin{problem}{题 4 对角线上的垂直与等长}
正方形 \(ABCD\) 中，\(E\) 为对角线 \(AC\) 内部的一点，\(N\) 在 \(CD\) 上，且 \(BE\perp EN\)。证明 \(BE=EN\)。
\diagram{\figFour}
\hint{连接 \(DE\)。关于对角线 \(AC\) 对称的两个顶点是哪两个？}
\end{problem}
\clearpage
\section{二 从三角形补出模型}
\begin{problem}{题 5 同一道题的三种构造}
\(\triangle ABC\) 中，\(D\) 为 \(BC\) 内部的一点，\(AD\perp BC\)，\(\angle BAC=45^\circ \)，\(BD=3\)，\(AD=6\)，求 \(CD\)。分别尝试旋转、两次翻折、延长构造等腰直角三角形。
\diagram{\figFive}
\hint{在直线 \(BC\) 上取 \(E\)、\(F\)，使 \(DE=DF=AD\)，且 \(E\)、\(B\)、\(D\)、\(C\)、\(F\) 按此顺序排列。}
\end{problem}
\section{三 矩形与交叉角}
\begin{problem}{题 6 矩形中截出正方形}
矩形 \(ABCD\) 中，\(AD=6\)，\(AB=12\)，\(F\in CD\)，\(DF=3\)，\(E\in BC\)，\(\angle EAF=45^\circ \)，求 \(BE\)。
\diagram{\figSix}
\hint{取 \(AB\) 的中点 \(M\)，过 \(M\) 作 \(MQ\parallel AD\)，\(Q\in DC\)。\(AE\) 与 \(MQ\) 交于 \(N\)。}
\end{problem}
\clearpage
\begin{problem}{题 7 把交叉角搬到顶点}
直角三角形 \(ABC\) 中，\(\angle C=90^\circ \)，\(AC=6\)，\(BC=8\)，\(D\) 是 \(AC\) 的中点，\(E\in BC\)，\(BD\) 与 \(AE\) 交于 \(F\)。若 \(\angle BFE=45^\circ \)，求 \(CE\)。
\diagram{\figSeven}
\hint{过 \(A\) 作 \(AG\perp BD\)，\(G\) 在 \(BC\) 向 \(C\) 外的延长线上，把交叉角转为 \(\angle GAE\)。}
\end{problem}
\begin{problem}{题 8 相等的两段补足旋转}
矩形 \(ABCD\) 中，\(E\)、\(F\) 分别在 \(BC\)、\(CD\) 的内部，\(CE=CF\)，\(DF=3\)，\(BE=4\)，\(\angle EAF=45^\circ \)，求 \(AB\)。
\diagram{\figEight}
\hint{把 \(AE\) 绕 \(A\) 逆时针旋转 \(90^\circ \)。旋转后的端点与 \(F\) 的水平距离、竖直距离能否直接算出？}
\end{problem}
\clearpage
\section{四 位置变化与分类讨论}
\begin{problem}{题 9 从线段和到线段差}
四边形 \(ABCD\) 的四条边相等，四个角都是直角，连接 \(AM\)、\(AN\)、\(MN\)。

\((1)M\)、\(N\) 分别在 \(BC\)、\(CD\) 的内部，\(\angle MAN=45^\circ \)。证明 \(MN=BM+DN\)。（原题 \(6\) 分）

\((2)M\) 在 \(BC\) 向 \(C\) 外的延长线上，\(N\) 在 \(CD\) 向 \(D\) 外的延长线上，仍有 \(\angle MAN=45^\circ \)。判断 \(BM\)、\(DN\)、\(MN\) 的数量关系，并证明。（原题 \(6\) 分）
\diagram{\figNine}
\hint{两问采用同一次旋转，再观察三个共线点的顺序。}
\end{problem}
\begin{problem}{题 10 一个条件对应两种构型}
\(\triangle ABC\) 中，\(D\) 在 \(BC\) 上，\(\angle ABC=45^\circ \)，\(\angle CAD=45^\circ \)，\(BD=4\)，\(BC=9\)，求 \(AB\) 的长。
\diagram{\figTen}
\hint{过 \(A\) 作 \(AH\perp BC\)，垂足 \(H\) 在直线 \(BC\) 上。先利用 \(\angle B=45^\circ \)，再讨论 \(H\) 与 \(D\) 的位置。}
\end{problem}
